\begin{filecontents}{bibliography.bib}
@mastersthesis{carneiro2019type,
  author     = {Mario Carneiro},
  title      = {The Type Theory of Lean},
  school     = {Carnegie Mellon University},
  year       = {2019},
  note       = {\url{https://github.com/digama0/lean-type-theory/releases/tag/v1.0}}
}
\end{filecontents}

\documentclass{article}

\usepackage[style=authoryear, backend=biber]{biblatex}
\addbibresource{bibliography.bib}

\usepackage[T1,T2A]{fontenc}
\usepackage[utf8]{inputenc}
\usepackage[english,russian]{babel}
\usepackage{indentfirst}

\usepackage[includehead]{geometry}
\pagenumbering{gobble}

\usepackage{fancyhdr}
\pagestyle{fancy}
\lhead{Алгебра, ИИ, Lean}
\chead{STLC и введение в Lean}
\rhead{22 октября 2026}

\usepackage{
    amssymb,    % \triangleq
    ebproof,    % proof trees
    hyperref,   % \href
    mathpartir, % inference rules
    simplebnf,  % BNF
    stmaryrd    % \llbracket, \rrbracket
}

\usepackage{amsthm}
\theoremstyle{definition}
\newtheorem{definition}{Определение}

\begin{document}

Lean --- одно из многих средств для интерактивных
доказательств. Все они устроены похожим образом: в
основе лежит некий минималистичный функциональный язык
программирования, программы на котором соответствуют
доказательствам в некоторой конструктивной логике; набор
некоторых внешних для этого языка средств упрощает
создание доказательств. В доказательствах с помощью LLM
чаще всего используют именно Lean по следующим причинам:

\begin{enumerate}
    \item Много уже существующего кода для обучения
        нейронных сетей (<<большой датасет>>);
    \item Использование принципов классической
        математики, упрощающее перевод;
    \item Крупнейшая библиотека формализованной
        математики
        \href{https://github.com/leanprover-community/mathlib4}{Mathlib},
        на которую можно ссылаться,
        не принимая сомнительных аксиом.
\end{enumerate}

Lean в полном объёме мы займёмся в следующий раз; здесь
мы рассмотрим основания языка программирования, лежащего
в его сердце, и поймём, почему программы на этом языке
вообще что-то доказывают.

\section{Просто типизированное лямбда-исчисление}

Здесь мы рассмотрим простейшее исчисление, в котором
наблюдается соответствие Карри-Говарда, согласно
которому программы оказываются доказательствами: просто
типизированное лямбда-исчисление (STLC, simply typed
lambda calculus). Чтобы произвести это соответствие, нам
нужно будет классифицировать доказательства по тому,
какое утверждение они доказывают; здесь нам пригодятся
\emph{типы}:

\begin{figure}[ht]
    \centering
    \begin{bnf}
        $A$ : \textsf{Атом} :in: $\mathcal{A}$;;
        $S$,$T$ : \textsf{Тип} ::=
        | $A$ : атомарный тип
        | $S\to T$ : тип-стрелка
        ;;
    \end{bnf}
    \caption{Грамматика типов STLC}
    \label{stlc-types-grammar}
\end{figure}

На Рис.~\ref{stlc-types-grammar} описана
\emph{грамматика} типов, то есть правила, по которым
можно получить любой тип системы. Нотация, в которой она
дана, известна как \emph{форма Бэкуса-Наура} и означает
следующее:

\begin{definition}[Типы STLC]
    Зафиксируем некоторый алфавит атомарных символов
    $\mathcal{A}$, элемент которого мы обычно будем
    обозначать $A$. \emph{Множеством типов} (элементы
    которого мы обычно будем обозначать $S$, $T$)
    является наименьшее множество синтаксических
    деревьев, удовлетворяющее следующим свойствам:
    \begin{itemize}
        \item Для всякого $A\in\mathcal{A}$ есть
            \emph{атомарный тип} $A$, являющийся типом;
        \item Если $S$ и $T$ --- типы, то $S\to T$ ---
            тоже тип, также называющийся
            <<стрелка из $S$ в $T$>>.
    \end{itemize}
\end{definition}

Обратите внимание, что к типам (и к программам, к
которым мы перейдём позже) правильнее относиться именно
как к синтаксическим деревьям, а не просто как к
строчкам. При записи таких деревьев в тексте, конечно,
нужно использовать скобки для однозначного разбора:
\[
    (S\to T)\to U \stackrel{?}{=} S\to T\to U
    \stackrel{?}{=} S\to(T\to U)
\]
Правда, для облегчения письма и чтения в таких случаях
принимают некоторые соглашения о нотации. Так, в случае
стрелки, $(\to)$ считается правоассоциативным
оператором, т.е. $S\to T\to U$ понимается как
$S\to(T \to U)$.

Множество програм, классифицируемое такими типами,
строится как на Рис.~\ref{stlc-terms-grammar}.

\begin{figure}[ht]
    \centering
    \begin{bnf}
        $x$ : \textsf{Переменная} :in: $\mathcal{V}$;;
        $s$,$t$ : \textsf{Терм} ::=
        | $x$ : переменная как терм
        | $\lambda x.t$ : анонимная функция
        | $t\,s$ : применение функции
        ;;
    \end{bnf}
    \caption{Грамматика термов STLC}
    \label{stlc-terms-grammar}
\end{figure}

Применение функции левоассоциативно: $(s\,t)\,u$
понимается как $s\,t\,u$. Кроме этого, при
последовательном объявлении анонимных функций все
символы $\lambda$ кроме первого можно опустить:
$\lambda x\,y\,z.\,t$ понимается как
$\lambda x.(\lambda y.(\lambda z.t))$.

Наконец, свяжем вместе термы и типы: скажем, какие
термы $t$ к какому типу $T$ относятся; этим занимается
\emph{суждение о типизации}, обозначаемое $t:T$.

Идея этого суждения состоит в том, что типы --- это
что-то вроде множеств, а термы типа $T$ --- элементы
соответствующего множества; тип-стрелка из $A$ в $B$
тогда не что иное, как множество функций из $A$ в $B$, а
$\lambda x.t$ --- форма записи функции, по иксу из $A$
возвращающей $t$ из $B$. Но тогда возникает вопрос ---
какой тип имеет переменная $x$? В суждении
$(\lambda x.x):T\to T$ она, очевидно, имеет тип $T$; но
если мы заменим $T$ на $U$, то икс уже будет иметь тип
$U$! По сути, тип икса зависит от \emph{контекста}:
списка, хранящего информацию о типе каждой переменной,
используемой в текущем выражении.

\begin{figure}[ht]
    \centering
    \begin{bnf}
        $\Gamma$ : \textsf{Контекст} ::=
        | $\cdot$ : пустой контекст
        | $\Gamma,x\,\colon\,T$ : расширение контекста
    \end{bnf}
    \caption{Контекст типизации STLC}
\end{figure}

Итого, суждение о типизации имеет вид $\Gamma\vdash t:T$
и читается как <<в контексте $\Gamma$, терм $t$ имеет
тип $T$>>. Задаётся оно с помощью \emph{правил вывода},
приведённых на Рис.~\ref{stlc-typing}.

\begin{figure}[ht]
    \begin{mathpar}
        \inferrule[Var]{ }
        {\Gamma\vdash x:T}
        (x:T)\in\Gamma
        \and
        \inferrule[Abs]
        {\Gamma,x:S\vdash t:T}
        {\Gamma\vdash\lambda x.t:S\to T}
        \and
        \inferrule[App]{
            \Gamma\vdash t:S\to T
            \\
            \Gamma\vdash s:S
        }{\Gamma\vdash t\,s:T}
    \end{mathpar}
    \caption{Правила типизации STLC}
    \label{stlc-typing}
\end{figure}

Опять же, правила вывода можно понимать просто как очень
краткую нотацию для следующего определения, по аналогии
с нотацией Бэкуса-Наура:

\begin{definition}[Отношение типизации в STLC]
    Отношение типизации $\Gamma\vdash t:T$ ---
    трёхместное отношение, связывающее вместе контекст
    $\Gamma$, терм $t$ и тип $T$, являющееся наименьшим
    отношением, удовлетворяющим следующим правилам:
    \begin{itemize}
        \item Если в контексте $\Gamma$ есть запись вида
            $x:T$, то $\Gamma\vdash x:T$.
        \item Если в расширенном контексте $\Gamma,x:S$
            терм $t$ имеет тип $T$, то
            $\Gamma\vdash\lambda x.t:S\to T$.
        \item Если в контексте $\Gamma$ верно, что
            $t:S\to T$ и $s:S$, то
            $\Gamma\vdash t\,s:T$.
    \end{itemize}
\end{definition}

В целом, уже построенную систему можно рассматривать как
способ формальной записи некоторых
теоретико-множественных функций. Например, мы можем
протипизировать композицию функций (ниже
$\Gamma \triangleq f:B\to C,g:A\to B,x:A$ для
краткости):

\begin{figure}[ht]
    \centering
    \begin{prooftree}
        \hypo{ }
        \infer1[\textsc{Var}]{\Gamma\vdash f:B\to C}
        \hypo{ }
        \infer1[\textsc{Var}]{\Gamma\vdash g:A\to B}
        \hypo{ }
        \infer1[\textsc{Var}]{\Gamma\vdash x:A}
        \infer2[\textsc{App}]{\Gamma\vdash g\,x:B}
        \infer2[\textsc{App}]{
            f:B\to C,g:A\to B,x:A\vdash f(g\,x):C
        }
        \infer1[\textsc{Abs}]{
            f:B\to C,g:A\to B\vdash\lambda x.f(g\,x):
            A\to C
        }
        \infer1[\textsc{Abs}]{
            f:B\to C\vdash\lambda g\,x.f(g\,x):
            (A\to B)\to A\to C
        }
        \infer1[\textsc{Abs}]{
            \cdot\vdash\lambda f\,g\,x.f(g\,x):
            (B\to C)\to(A\to B)\to A\to C
        }
    \end{prooftree}
    \caption{Вывод типизации композиции функций}
\end{figure}

Теперь давайте сотрём все термы из вывода выше.
Получим следующее дерево вывода в \emph{импликативном
фрагменте интуиционистской пропозициональной логики}:

\begin{figure}[ht]
    \centering
    \begin{prooftree}
        \hypo{ }
        \infer1{\Gamma\vdash B\to C}
        \hypo{ }
        \infer1{\Gamma\vdash A\to B}
        \hypo{ }
        \infer1{\Gamma\vdash A}
        \infer2{\Gamma\vdash B}
        \infer2{B\to C,A\to B,A\vdash C}
        \infer1{B\to C,A\to B\vdash A\to C}
        \infer1{B\to C\vdash(A\to B)\to A\to C}
        \infer1{
            \cdot\vdash(B\to C)\to(A\to B)\to A\to C
        }
    \end{prooftree}
    \caption{Вывод в Int}
\end{figure}

Это конкретный пример применения \emph{соответствия
Карри-Говарда}: любое суждение о типизации
$\Gamma\vdash t:T$ соответствует доказательству
$(\downarrow\Gamma)\vdash T$ \footnote{Где
$\downarrow\cdot=\cdot$ и $\downarrow(\Gamma,x:T)=
(\downarrow\Gamma),T$.}; и любое доказательство
$\Gamma_0\vdash T$ соответствует суждению о типизации
некоторого терма $(\uparrow\Gamma_0)\vdash t:T$
\footnote{Где $\uparrow\cdot=\cdot$ и
$\uparrow(\Gamma,T)=(\uparrow\Gamma),x_T:T$ с
точностью до выбора новых имён для переменных
одинакового типа.}.

В чём плюс от введения нового исчисления? Во-первых, у
нас теперь есть единый формализм, в рамках которого
можно \emph{и} строить теоретико-множественные объекты
(где $A\to B$ интерпретируется как множество функций),
\emph{и} доказывать какие-то утверждения (где $A\to B$
интерпретируется как импликация). Во-вторых, задача
проверки типизации в STLC решается алгоритмически ---
другими словами, по $\Gamma$, $t$ и $T$ можно проверить
автоматически, верно ли, что $\Gamma\vdash t:T$. Так что
пользователю достаточно выписать тип и программу, а
корректность логического вывода среда интерактивных
доказательств проверит сама.

\section{Вновь о Lean}

Конечно, в Lean используется чуть более сложная система,
чем STLC: для содержательных математических рассуждений
нам не хватает кванторов, других логических связок кроме
импликации, способов строить бесконечные множества,
фактормножества и так далее. Но, надеюсь, из изложения
выше понятен общий принцип работы: пользователь вместе с
интерактивной средой доказательств любым доступным ему
образом строит некоторые $t$ и $T$; ядро среды
проверяет, что $t:T$. По теоретико-множественной
семантике $\llbracket\_\rrbracket$,
доказанной\footnote{\fullcite{carneiro2019type}} для
Lean, полученной программе соответствует некоторый
теоретико-множественный объект
$\llbracket t\rrbracket\in\llbracket T\rrbracket$. Если
при этом $T$ было типом, соответствующим утверждению, то
$\llbracket t\rrbracket$ оказывается доказательством
этого утверждения.

Для практики с ядром Lean я составил несколько
упражнений, доступных в файле \\
\texttt{AiLeanMM/Basic.lean} в репозитории, в который я
буду выкладывать все свои материалы по данному
мини-курсу: \url{https://github.com/TurtlePU/AiLeanMM}.
Для удобства работы советую настроить Lean в
каком-нибудь редакторе кода; проще всего это сделать для
VS Code:
\begin{enumerate}
    \item Установите \texttt{VS Code}.
    \item Установите расширение \texttt{Lean 4}
        для VS Code.
    \item Склонируйте репозиторий \texttt{AiLeanMM}
        через стартовое меню VS Code.
    \item Следуйте указаниям расширения Lean 4
        (возможно, нужно сделать что-то с Вашим
        подключением к интернету, чтобы расширение
        смогло установить все необходимые пакеты на Ваш
        компьютер).
\end{enumerate}

\end{document}
